\chapter{Network-on-Chip (NoC)}\label{ch:hw-noc}

\section{Overview}
The blocks of Libre-V are tied together by an AXI4 Network-on-Chip generated with
\href{https://github.com/pulp-platform/FlooNoC}{FlooNoC}. Every off-subsystem transfer --- the Pico
reaching a peripheral, the Rocket fetching from its SRAMs, the DMA moving data, an interrupt-driving
device being programmed --- crosses the NoC. It is described declaratively in \texttt{noc\_desc.yml} and
elaborated by \texttt{floogen}; the base configuration uses a single router.

Endpoints are either \emph{managers} (they issue transactions) or \emph{subordinates} (they are
addressed). Managers carry no address range; subordinates own a fixed window in the 32-bit space.

\section{Manager Endpoints}

\begin{tabular}{@{}ll@{}}
\toprule
\textbf{Endpoint} & \textbf{Master} \\
\midrule
\texttt{pico\_subsystem}        & Pico core / COMMCTRL, via the AHB-to-AXI bridge (\autoref{ch:hw-pico-sys}). \\
\texttt{rocket\_subsystem\_mem} & Rocket cacheable port (\texttt{mem\_axi4}). \\
\texttt{rocket\_subsystem\_mmio}& Rocket uncacheable port (\texttt{mmio\_axi4}). \\
\texttt{dma\_mngr}              & DMA datapath (the read/write bursts). \\
\bottomrule
\end{tabular}

\section{Subordinate Endpoints and Address Map}
The subordinate windows below form the system address map for everything reached over the NoC. Local
Pico subordinates (\autoref{ch:sw-pico-sys}) and the Rocket-private devices (\autoref{ch:sw-rocket-sys})
do not appear here --- they never leave their subsystem.

\begin{longtable}{@{}llll@{}}
\toprule
\textbf{Endpoint} & \textbf{Base} & \textbf{Size} & \textbf{Device} \\
\midrule
\endhead
\texttt{uart1}            & \texttt{0x2000\_0000} & 64\,KB  & UART1 (\autoref{ch:sw-uart}) \\
\texttt{timer1}           & \texttt{0x2010\_0000} & 64\,KB  & TIMER1 (\autoref{ch:sw-timer}) \\
\texttt{dma\_sub}         & \texttt{0x2020\_0000} & 64\,KB  & DMA registers (\autoref{ch:sw-dma}) \\
\texttt{accelerator}      & \texttt{0x3000\_0000} & 32\,MB  & Accelerator (\autoref{ch:sw-accel}) \\
\texttt{sram\_imem\_rocket}& \texttt{0x8000\_0000} & 256\,KB & Rocket instruction SRAM \\
\texttt{sram\_dmem\_rocket}& \texttt{0x8010\_0000} & 256\,KB & Rocket data SRAM \\
\texttt{hyperram}         & \texttt{0xA000\_0000} & 256\,MB & External HyperRAM data window \\
\bottomrule
\end{longtable}

\begin{docwarning}
The peripheral band starts at \texttt{0x2000\_0000}, not \texttt{0x1000\_0000}: COMMCTRL intercepts
\texttt{0x1000\_0000} (bus handover) and \texttt{0x1000\_0004} (host interrupt) in its own front-end
before those addresses reach the bus (\autoref{ch:sw-pico-sys}), so \texttt{0x1000\_0000}--\texttt{0x1000\_0FFF}
is reserved and must stay unmapped. Peripherals are spaced 1\,MB apart within the band.
\end{docwarning}

\section{Clock Domains}
Each endpoint belongs to a clock domain, and the top level places an AXI clock-domain-crossing
(\texttt{axi\_cdc}) on the NoC side of each so a block can run at its own frequency. The base
configuration exposes five domains (Pico, Rocket, and one each broadly for the peripheral, HyperRAM, and
accelerator groups), clocked and reset through the CRG (\autoref{ch:sw-crg}).

\begin{docnote}
In the current simulation/testbench setup every domain is driven by the same clock, so the crossings
operate as same-clock passthroughs. The CRG's per-domain clock generation and true frequency ratios are
not exercised in that mode.
\end{docnote}

\section{Top-Level Integration}
\texttt{libre\_top} instantiates the NoC (\texttt{floo\_libre\_v\_noc}) together with every endpoint's
NoC-side wrapper and the clock-domain crossings:

\begin{tabular}{@{}ll@{}}
\toprule
\textbf{Instance} & \textbf{Role} \\
\midrule
\texttt{pico\_subsystem}        & Control plane + Pico manager (\autoref{ch:hw-pico-sys}) \\
\texttt{rocket\_subsystem}      & Rocket manager, cacheable + uncacheable ports (\autoref{ch:hw-rocket-sys}) \\
\texttt{noc\_dma}               & DMA manager + register subordinate \\
\texttt{axi\_mem\_16384x64} ($\times$2) & Rocket instruction/data SRAMs \\
\texttt{noc\_accelerator}       & Accelerator subordinate \\
\texttt{noc\_hyperram}          & HyperRAM controller + data window \\
\texttt{noc\_axi\_uart}         & UART1 subordinate \\
\texttt{noc\_axi\_timer}        & TIMER1 subordinate \\
\texttt{floo\_libre\_v\_noc}    & The generated FlooNoC fabric (router + links) \\
\texttt{axi\_cdc} ($\times$N)   & Per-endpoint clock-domain crossings \\
\bottomrule
\end{tabular}

\section{Source Files}

\begin{tabular}{@{}ll@{}}
\toprule
\textbf{File} & \textbf{Role} \\
\midrule
\texttt{top/libre\_top.sv} & SoC top-level: NoC + all endpoints + clock-domain crossings \\
\texttt{noc/noc\_desc.yml} & Declarative NoC description (endpoints, address map, router) \\
\texttt{noc/hw/} & Generated FlooNoC RTL \\
\bottomrule
\end{tabular}

\newpage
